Знаменитый чтец Антон Исаакович Ш. "--- то самое историческое лицо, которое выступало в сентябре месяце 1940 года в Литейном лектории, "--- любило перед своими концертами полежать часок-другой и отдохнуть. Ляжет оно, бывало, на кушет и скажет:

<<Буду спать>>, "--- а само не спит.

После концертов оно любило поужинать.

Вот оно придет домой, рассядется за столом и говорит своей жене:

<<А ну, голубушка, состряпай"=ка мне что"=нибудь из лапши.>>

И пока жена его стряпает, оно сидит за столом и книгу читает.

Жена его хорошенькая, в кружевном передничке, с сумочкой в руках, а в сумочке кружевной платочек и ватрушечный медальончик лежат, жена его бегает по комнате, каблучками стучит, как бабочки, а оно скромно за столом сидит, ужина дожидается.

Все так складно и прилично. Жена ему что"=нибудь приятное скажет, а оно головой кивает. А жена порх к буфетику и уже рюмочками там звенит.

<<Налей"=ка, душечка, мне рюмочку>>, "--- говорит оно.

<<Смотри, голубчик, не спейся>>, "--- говорит ему жена.

<<Авось, пупочка, не сопьюсь>>, "--- говорит оно, опрокидывая рюмочку в рот.

А жена грозит ему пальчиком, а сама боком через двери на кухню бежит.

Вот в таких приятных тонах весь ужин проходит. А потом они спать закладываются.

Ночью, если им мухи не мешают, они спят спокойно, потому что уж очень они люди хорошие!

\medskip
\noindent\textit{1940 г.}
